\section{Listening}

\subsection{Unit 1}

\subsubsection{Conversation}
\begin{enumerate}
	\item What does the woman normally do in the morning?
	\begin{choices}
		\item Get up late.
		\item Drink a lot of coffee.
		\item Make her work plans.
		\item \textbf{\color{red}Do physical activity.}
	\end{choices}

	\item What is one advantage of the man's schedule?
	\begin{choices}
		\item He can arrange more time for social life at night.
		\item He can have less difficulty coping with meetings.
		\item \textbf{\color{red}He can enjoy a less distracting work environment.}
		\item He can be more focused and productive during the day.
	\end{choices}

	\item What do the speakers think determines whether someone is an early bird or a night owl?
	\begin{choices}
		\item One's sleep quality.
		\item One's work schedule.
		\item \textbf{\color{red}One's individual body clock.}
		\item One's attitude toward morning exercise.
	\end{choices}

	\item What should individuals do regarding their schedules according to the speakers?
	\begin{choices}
		\item Sleep less to get more work done.
		\item Change their natural rhythm to match that of others.
		\item Follow a fixed routine regardless of external distractions.
		\item \textbf{\color{red}Listen to their body and adjust their schedule accordingly.}
	\end{choices}
\end{enumerate}

\subsubsection{Passage}
\begin{enumerate}
	\item What tasks does the speaker tend to do first every day?
	\begin{choices}
		\item \textbf{\color{red}The quick tasks.}
		\item The routine tasks.
		\item The most urgent tasks.
		\item The most difficult tasks.
	\end{choices}

	\item What is one of the suggestions for getting one's work done in time?
	\begin{choices}
		\item Getting up as early as possible to work.
		\item Being a night owl and working late at night.
		\item Using emails to improve work efficiency.
		\item \textbf{\color{red}Making use of one's most productive time of day.}
	\end{choices}

	\item How does the speaker deal with emails now?
	\begin{choices}
		\item She checks and responds to them constantly.
		\item She accesses them only during her real work hours.
		\item She no longer responds to the most trivial of them.
		\item \textbf{\color{red}She opens her inbox only when it's totally necessary.}
	\end{choices}

	\item What is the purpose of the talk?
	\begin{choices}
		\item To discuss the importance of email.
		\item To show how to balance work and life.
		\item To share some ideas on how to reduce workload.
		\item \textbf{\color{red}To give some tips on how to work more efficiently.}
	\end{choices}
\end{enumerate}

\subsubsection{Lecture 1}
\begin{enumerate}
	\item What can we learn about verbal behaviors from the lecture?
	\begin{choices}
		\item \textbf{\color{red}They include both spoken and written language.}
		\item They ensure that what we say is what we mean.
		\item They can convey messages about our personality.
		\item They are more complex than nonverbal behaviors.
	\end{choices}

	\item What does paralanguage refer to?
	\begin{choices}
		\item Expressing emotions through writing.
		\item Making requests with verbal language.
		\item Using body language to convey meaning.
		\item \textbf{\color{red}Using voice in particular ways to communicate.}
	\end{choices}

	\item What does the speaker mean by saying "behavior doesn't occur in isolation"?
	\begin{choices}
		\item Behavior makes communication complex.
		\item Behavior may cause offense to other people.
		\item \textbf{\color{red}Behavior can be influenced by cultural factors.}
		\item Behavior plays a vital role in conveying meaning.
	\end{choices}

	\item What is the main purpose of the lecture?
	\begin{choices}
		\item To define verbal and nonverbal behaviors.
		\item \textbf{\color{red}To illustrate the communicative function of behavior.}
		\item To explain how context affects the meaning of our words.
		\item To show how to behave appropriately in different cultures.
	\end{choices}
\end{enumerate}

\subsubsection{Lecture 2}
\begin{enumerate}
	\item What does the American psychologist suggest as the reason for people failing to implement new behaviors consistently?
	\begin{choices}
		\item \textbf{\color{red}Lack of a proper approach.}
		\item Lack of motivation or willpower.
		\item The disruption of unhealthy habits.
		\item Frequent adjustments of the desired goal.
	\end{choices}

	\item What is recommended by the psychologist as a way to lose weight and improve physical health?
	\begin{choices}
		\item Building a better bedtime routine.
		\item Following a strict diet and exercise routine.
		\item Eating less and doing exercise in the morning.
		\item \textbf{\color{red}Starting by making small changes like taking short walks.}
	\end{choices}

	\item Which is mentioned as a small step for someone to achieve a career goal?
	\begin{choices}
		\item Taking the time to actively listen to others.
		\item Being conscious enough to express gratitude.
		\item \textbf{\color{red}Seeking out mentors to gain guidance and support.}
		\item Building stronger working relationships with others.
	\end{choices}

	\item Why can the "small steps" approach lead to lasting change?
	\begin{choices}
		\item It can avoid exhaustion and low motivation.
		\item It guarantees immediate results and success.
		\item \textbf{\color{red}It breaks down goals into manageable tasks.}
		\item It reduces the time and effort for achieving an objective.
	\end{choices}
\end{enumerate}

\newpage
\subsection{Unit 2}

\subsubsection{Conversation}
\begin{enumerate}
	\item What is one of the things the woman wishes for at first?
	\begin{choices}
		\item Learning to swim.
		\item \textbf{\color{red}Finding a well-paid job.}
		\item Seeing her husband soon.
		\item Getting married in a big house.
	\end{choices}

	\item What difficulties did the man experience in his life?
	\begin{choices}
		\item He had a hard time keeping his marriage.
		\item He got really sick and lost all his valuable things.
		\item \textbf{\color{red}He lost his job and wealth and was sick for years.}
		\item He was unable to see his children due to poor health.
	\end{choices}

	\item According to the man, why is true love important for a happy life?
	\begin{choices}
		\item It helps one value little things in life.
		\item \textbf{\color{red}It makes couples appreciate each other.}
		\item It gives people patience in their job search.
		\item It makes people confident in pursuing their goals.
	\end{choices}

	\item What is the purpose of the man striking up the conversation?
	\begin{choices}
		\item To tell the woman an interesting life story.
		\item To teach the woman skills to find a good job.
		\item To encourage the woman to pursue a happy marriage.
		\item \textbf{\color{red}To give the woman advice on the important things in life.}
	\end{choices}
\end{enumerate}

\subsubsection{Passage}
\begin{enumerate}
	\item What is predicted about the human population in the next three decades?
	\begin{choices}
		\item The proportion of 15- to 50-year-olds will decrease slightly.
		\item The number of people aged 80 and above will increase by about 22\%.
		\item \textbf{\color{red}The proportion of people aged 60 and above will increase significantly.}
		\item The number of 12- to 22-year-olds will decrease by more than 420 million.
	\end{choices}

	\item What is one impact of an aging human population on society?
	\begin{choices}
		\item Economic productivity is constantly rising.
		\item The range of healthcare services is expanding.
		\item The pension system will be replaced in the near future.
		\item \textbf{\color{red}The ratio of working-age people to older people is decreasing.}
	\end{choices}

	\item What may be a result of the population changes discussed in the passage?
	\begin{choices}
		\item Instability of housing prices.
		\item \textbf{\color{red}Changes in family structure.}
		\item An increase in property taxation.
		\item A decrease in population mobility.
	\end{choices}

	\item What does human society need according to the speaker?
	\begin{choices}
		\item Strong trade ties between countries around the world.
		\item Effective policies to reduce social and health costs.
		\item \textbf{\color{red}Necessary changes to address the needs of older people.}
		\item Rich resources to maintain the traditional family structure.
	\end{choices}
\end{enumerate}

\subsubsection{Lecture 1}
\begin{enumerate}
	\item What is the most reliable predictor of one's chances of getting age-related diseases?
	\begin{choices}
		\item Year of birth.
		\item \textbf{\color{red}Biological age.}
		\item Chronological age.
		\item Number of body cells.
	\end{choices}

	\item How is biological age measured?
	\begin{choices}
		\item By comparing how old one looks and how old one really is.
		\item By identifying the diseases one has suffered in chronological order.
		\item \textbf{\color{red}By examining biological evidence of one's physical and mental functions.}
		\item By analyzing the current health status of one's family members and relatives.
	\end{choices}

	\item What does the speaker say about the importance of genetic influence on biological age?
	\begin{choices}
		\item It has not been studied until recent decades.
		\item It cannot be measured effectively through DNA research.
		\item \textbf{\color{red}It has been supported by a lot of research in recent decades.}
		\item It is equal to that of environmental factors and lifestyle choices.
	\end{choices}

	\item What does the speaker say about healthy lifestyle choices?
	\begin{choices}
		\item \textbf{\color{red}They can help reverse one's biological age.}
		\item They can improve one's exercise routine in the long run.
		\item They can reduce one's exposure to stressful environments.
		\item They can benefit the economic status of certain population groups.
	\end{choices}
\end{enumerate}

\subsubsection{Lecture 2}
\begin{enumerate}
	\item What does the speaker mainly talk about?
	\begin{choices}
		\item Distinctions between concrete thinking and abstract thinking.
		\item Similarities between cognitive skills and formal logical thinking.
		\item \textbf{\color{red}Features of cognitive development in different phases of adolescence.}
		\item Definitions of three thinking patterns developed during the teenage period.
	\end{choices}

	\item What does complex thinking focus on during one's early adolescent years?
	\begin{choices}
		\item \textbf{\color{red}Decisions about personal life in school and home settings.}
		\item Development of a code of ethics and personal relationships.
		\item Arrangement and order of concrete objects and simple actions.
		\item Analysis and understanding of other people's ideas and opinions.
	\end{choices}

	\item During which period are children most likely to think a lot about global issues?
	\begin{choices}
		\item Ages 6-12.
		\item Ages 13-15.
		\item \textbf{\color{red}Late adolescence.}
		\item Early adolescence.
	\end{choices}

	\item What does the speaker suggest parents and teachers should do?
	\begin{choices}
		\item Encourage children to focus on their current goals.
		\item Be patient with children's mistakes in cognitive thinking.
		\item Give children the freedom to keep their thoughts to themselves.
		\item \textbf{\color{red}Support the healthy development of cognitive skills in children.}
	\end{choices}
\end{enumerate}

\newpage
\subsection{Unit 3}

\subsubsection{Conversation}
\begin{enumerate}
	\item What is one reason that makes Big Sky Park attractive to the couple?
	\begin{choices}
		\item It is right in their city.
		\item It is an entertainment park.
		\item \textbf{\color{red}It has nice trails and camping sites.}
		\item It's a park with a beautiful city skyline.
	\end{choices}

	\item What problem do they try to figure out about going to the park?
	\begin{choices}
		\item Where the No. 32 bus stop is.
		\item What the best route to drive to the park is.
		\item Whether they can find easy parking for their car.
		\item \textbf{\color{red}How they can get their camping equipment on the bus.}
	\end{choices}

	\item What will the woman do before their trip?
	\begin{choices}
		\item Buy a big backpack.
		\item \textbf{\color{red}Bake a chocolate cake.}
		\item Order their favorite food.
		\item Make a detailed trip plan.
	\end{choices}

	\item What does the couple long to do at the park?
	\begin{choices}
		\item Drink coffee around a campfire.
		\item Have a leisurely barbecue dinner.
		\item \textbf{\color{red}Walk together in the beauty of nature.}
		\item Take a bike ride around the open spaces.
	\end{choices}
\end{enumerate}

\subsubsection{Passage}
\begin{enumerate}
	\item Why should visitors hurry if they want to see the exhibits in the new museum?
	\begin{choices}
		\item \textbf{\color{red}The exhibition will last for no more than a week.}
		\item The museum opens only during the holiday season.
		\item The museum workers will soon begin their vacation.
		\item Large crowds of visitors will come to the unusual show.
	\end{choices}

	\item What is exhibited in the new museum?
	\begin{choices}
		\item Plenty of televisions that look stylish.
		\item Anything that reflects modern fashion.
		\item \textbf{\color{red}Beds and sofas that make people feel relaxed.}
		\item Objects that have brought big lifestyle changes.
	\end{choices}

	\item What is the purpose of the exhibition?
	\begin{choices}
		\item To offer a place to rest in our fast-paced societies.
		\item \textbf{\color{red}To provide a practical experience for people to think about laziness.}
		\item To teach lazy people how to change their behavior and lifestyle.
		\item To show how furniture has changed to meet the needs of modern people.
	\end{choices}

	\item What did the museum's founder advise us to do?
	\begin{choices}
		\item Consider laziness an enemy of work.
		\item \textbf{\color{red}Slow down, relax, and be lazy for a while.}
		\item Make efforts to shorten our working week.
		\item Enjoy our leisure time if our jobs allow us to.
	\end{choices}
\end{enumerate}

\subsubsection{Lecture 1}
\begin{enumerate}
	\item What is the main feature of niche tourism?
	\begin{choices}
		\item It offers budget-friendly travel packages.
		\item It focuses on popular tourist destinations.
		\item \textbf{\color{red}It provides opportunities for unique experiences.}
		\item It specializes in historical and geographical sites.
	\end{choices}

	\item Which types of niche travel experienced the biggest growth from 2018 to 2021?
	\begin{choices}
		\item Adventure and wildlife trips.
		\item Wellness and friendship trips.
		\item Low-risk and conventional trips.
		\item \textbf{\color{red}Cookery and eco-friendliness trips.}
	\end{choices}

	\item What tip is offered for finding niche travel?
	\begin{choices}
		\item \textbf{\color{red}Consulting with a travel advisor.}
		\item Following the advice of tourism analysts.
		\item Doing thorough research through various sources.
		\item Contacting an operator focusing on immersive vacations.
	\end{choices}

	\item What does the lecture suggest about the future of niche travel?
	\begin{choices}
		\item It will decline in popularity.
		\item It will attract all age groups.
		\item It will remain a niche market.
		\item \textbf{\color{red}It will become fairly mainstream.}
	\end{choices}
\end{enumerate}

\subsubsection{Lecture 2}
\begin{enumerate}
	\item What can we learn about the development of the term "forest bathing" from the lecture?
	\begin{choices}
		\item \textbf{\color{red}It first originated in Japan in the early 1980s.}
		\item It was initially used to refer to swimming in forest ponds.
		\item It was first used by people who did regular walking in nature.
		\item It was coined by scientists studying the effects of nature on health.
	\end{choices}

	\item How long does a person need to spend in nature each week to experience the benefits of forest bathing, according to one study?
	\begin{choices}
		\item No more than 10 minutes.
		\item No more than 20 minutes.
		\item \textbf{\color{red}At least 120 minutes.}
		\item At least 250 minutes.
	\end{choices}

	\item Besides the duration of time spent in nature, what is another factor that may contribute to the positive effects of forest bathing?
	\begin{choices}
		\item Being exposed to new experiences.
		\item \textbf{\color{red}Experiencing feelings of awe in nature.}
		\item Resolving personal problems and concerns.
		\item Engaging in high-intensity physical exercise.
	\end{choices}

	\item Which of the following statements is supported by scientific literature?
	\begin{choices}
		\item \textbf{\color{red}Forest bathing has powerful restorative effects.}
		\item Forest bathing provides rather short-term benefits.
		\item Forest bathing can be regarded as a kind of medical treatment.
		\item Forest bathing becomes increasingly effective as time spent in nature increases.
	\end{choices}
\end{enumerate}

\newpage
\subsection{Unit 4}

\subsubsection{Conversation}
\begin{enumerate}
	\item What are the speakers talking about?
	\begin{choices}
		\item \textbf{\color{red}Potential choices for an art feature for their town.}
		\item Some local cinemas hosting entertainment activities.
		\item Job opportunities created by some new art projects in their area.
		\item Proven benefits of having diverse art activities in the community.
	\end{choices}

	\item What does the woman think of the cinema project?
	\begin{choices}
		\item It's her favorite because of its state-of-the-art technology.
		\item It's her top choice because of its economic benefits for the town.
		\item It's not an ideal project due to possibly limited use by young people.
		\item \textbf{\color{red}It's not the best choice for the town due to the high construction costs.}
	\end{choices}

	\item What does the woman say about the botanical garden project?
	\begin{choices}
		\item It lacks an appealing variety of plants and birds.
		\item It has a poor environment for immersive activities.
		\item \textbf{\color{red}It allows visitors to interact and connect with nature.}
		\item It provides great educational opportunities for adults.
	\end{choices}

	\item What will help lower the maintenance costs of the botanical garden?
	\begin{choices}
		\item \textbf{\color{red}An irrigation system with the latest water conservation technology.}
		\item An eco-friendly color for facilities to match the natural environment.
		\item A volunteer group that builds a durable and self-sustaining irrigation system.
		\item A building design that adopts the strictest maintenance standards for gardens.
	\end{choices}
\end{enumerate}

\subsubsection{Passage}
\begin{enumerate}
	\item What does the passage mainly discuss?
	\begin{choices}
		\item The process of painting the Mona Lisa.
		\item \textbf{\color{red}The reasons for the fame of the Mona Lisa.}
		\item The influence of the Mona Lisa on Italian art.
		\item The techniques used in preserving the Mona Lisa.
	\end{choices}

	\item Who is the person painted in the Mona Lisa?
	\begin{choices}
		\item A woman who liked painting.
		\item A woman who was a merchant.
		\item \textbf{\color{red}The wife of a wealthy merchant.}
		\item The wife of a famous Italian artist.
	\end{choices}

	\item What made it possible for the Mona Lisa to survive for centuries?
	\begin{choices}
		\item The quality of the oil paints on the surface.
		\item The texture of the frame holding the painting.
		\item The type of canvas used to protect the wood frame.
		\item \textbf{\color{red}The use of a wood panel as the material to paint on.}
	\end{choices}

	\item What do we learn about the painting style and techniques used in the Mona Lisa?
	\begin{choices}
		\item The landscape was painted with a unique close-up view.
		\item The subject's eyes and nose were painted with shadows.
		\item Some brush marks were left in the background on purpose.
		\item \textbf{\color{red}The subject's hands and face were painted with great precision.}
	\end{choices}
\end{enumerate}

\subsubsection{Lecture 1}
\begin{enumerate}
	\item What is the main characteristic of representational art?
	\begin{choices}
		\item It portrays clear and detailed views of the land.
		\item \textbf{\color{red}It depicts objects with great similarity to things in real life.}
		\item It presents mythological scenes in historical real-world settings.
		\item It avoids the representation of religious matters in sculptures and paintings.
	\end{choices}

	\item What does the speaker say about the styles of representational art?
	\begin{choices}
		\item \textbf{\color{red}They have different degrees of abstraction now.}
		\item They were once little known by the general public.
		\item They have gradually lost popularity in the modern era.
		\item They do not favor abstract depiction of actual objects.
	\end{choices}

	\item What is a distinctive feature of non-representational art?
	\begin{choices}
		\item Simple shapes are often combined in a symmetrical way.
		\item Artists use multiple lines and colors to depict real objects.
		\item \textbf{\color{red}Viewers can have different interpretations of a piece of work.}
		\item Abstract objects are portrayed by the use of thick, straight lines.
	\end{choices}

	\item What can be inferred from the lecture?
	\begin{choices}
		\item How line patterns are used in a painting shows an artist's style.
		\item What color a painting focuses on can have rich abstract meanings.
		\item \textbf{\color{red}People may sometimes disagree on the genre of a particular painting.}
		\item The meaning of an artwork is determined by the genre under which it falls.
	\end{choices}
\end{enumerate}

\subsubsection{Lecture 2}
\begin{enumerate}
	\item What does the speaker say about John Denver's music?
	\begin{choices}
		\item It was little known outside his home country.
		\item It was popular for a very short period of time.
		\item \textbf{\color{red}It has received a large number of prizes and honors.}
		\item It reflects the musical trends of the past four decades.
	\end{choices}

	\item Why did John Denver choose to use "Denver" as part of his stage name?
	\begin{choices}
		\item Denver was the name of his best friend in his earlier life.
		\item \textbf{\color{red}Denver is the name of the capital city of his favorite state.}
		\item He wanted to get away from the music scene of Los Angeles.
		\item He wanted to seek inspiration from the scenery of the Rocky Mountains.
	\end{choices}

	\item What was the first important event that led to John Denver's fame in the music industry?
	\begin{choices}
		\item He broke up with a band to perform his own songs.
		\item He joined a band to gain experience in songwriting.
		\item \textbf{\color{red}He was selected as lead singer for a popular folk band.}
		\item He was asked to produce a solo album in his own style.
	\end{choices}

	\item What contributed to John Denver's success in gaining a wide audience?
	\begin{choices}
		\item \textbf{\color{red}The use of elements from various musical genres.}
		\item The emphasis on authentic folk music traditions.
		\item Deep emotional connections with nature-loving audiences.
		\item Simple short lyrics related to themes of goodness and love.
	\end{choices}
\end{enumerate}

\newpage
\subsection{Unit 5}

\subsubsection{Conversation}
\begin{enumerate}
	\item Why does the man look upset?
	\begin{choices}
		\item He isn't satisfied with his team.
		\item \textbf{\color{red}He has too much work to complete.}
		\item He has met someone he doesn't like.
		\item He is afraid of his demanding manager.
	\end{choices}

	\item How does the woman say she can help the man?
	\begin{choices}
		\item She can get more people to work on his team.
		\item She can join his team to work for him herself.
		\item \textbf{\color{red}She can help him do his work in her free time.}
		\item She can get his manager to reduce his workload.
	\end{choices}

	\item What advice does the woman give to the man?
	\begin{choices}
		\item He should learn to deal with pressure.
		\item \textbf{\color{red}He should have confidence in himself.}
		\item He should feel free to ask others for help.
		\item He should accept more people into his team.
	\end{choices}

	\item Why is the restaurant called 365?
	\begin{choices}
		\item It serves 365 different dishes.
		\item It is open for business every day.
		\item It is located at No. 365 on Sixth Street.
		\item \textbf{\color{red}It offers seasonal foods all year round.}
	\end{choices}
\end{enumerate}

\subsubsection{Passage}
\begin{enumerate}
	\item What prediction is made by the global organization mentioned in the passage?
	\begin{choices}
		\item Type II diabetes will become the most common kind of disease.
		\item Diabetes rates among children and young people will be increasing.
		\item The issue of being overweight can hardly be solved in many countries.
		\item \textbf{\color{red}Over half of the world's population could be overweight or obese by 2035.}
	\end{choices}

	\item What are considered the primary causes of the growing issue of being overweight or obese?
	\begin{choices}
		\item Genetic factors and the increase in age.
		\item Increased risks of disease and stress in life.
		\item \textbf{\color{red}High-calorie foods and insufficient exercise.}
		\item High blood sugar and poor body composition.
	\end{choices}

	\item Compared to men who were just overweight, how much higher is the risk of diabetes for men who were obese, according to one study?
	\begin{choices}
		\item Three times higher.
		\item Four times higher.
		\item \textbf{\color{red}Five and a half times higher.}
		\item Three and a half times higher.
	\end{choices}

	\item What is the primary purpose of the passage?
	\begin{choices}
		\item To give suggestions on how to maintain a healthy body weight.
		\item To present statistics on the global issue of overweight and obesity.
		\item \textbf{\color{red}To explain the link between weight and the risk of diabetes in older adults.}
		\item To report how researchers studied the risk of diabetes in older people.
	\end{choices}
\end{enumerate}

\subsubsection{Lecture 1}
\begin{enumerate}
	\item What did the researchers of the new study find about the meaning of happiness?
	\begin{choices}
		\item \textbf{\color{red}It varies depending on geographical contexts.}
		\item It is influenced by a self-centered worldview.
		\item It can be strongly affected by people's lifestyle.
		\item It is primarily determined by personal success.
	\end{choices}

	\item How have Eastern ideologies influenced the concept of happiness in East Asian cultures?
	\begin{choices}
		\item Happiness is conditional and varies from person to person.
		\item Happiness is a process that connects everyone and everything.
		\item \textbf{\color{red}Happiness is closely linked to the quality of social relationships.}
		\item Happiness is defined by one's adherence to traditional values.
	\end{choices}

	\item In which countries did the happiness scale that emphasized independence prove to be most reliable according to the study?
	\begin{choices}
		\item In Asian countries.
		\item In African countries.
		\item In Central American countries.
		\item \textbf{\color{red}In Western European countries.}
	\end{choices}

	\item What does the speaker intend to discuss in the lecture?
	\begin{choices}
		\item The measurements of happiness across different countries.
		\item \textbf{\color{red}The influence of cultural values on the concept of happiness.}
		\item The difference between the Western and Eastern happiness scales.
		\item The potential negative effects of a universal standard of happiness.
	\end{choices}
\end{enumerate}

\subsubsection{Lecture 2}
\begin{enumerate}
	\item How do we typically set conditions for our happiness according to the speaker?
	\begin{choices}
		\item By constantly changing our life goals.
		\item \textbf{\color{red}By waiting for the perfect time to be happy.}
		\item By setting small goals for our everyday life.
		\item By comparing our conditions to those of others.
	\end{choices}

	\item Why does the speaker mention the study conducted by Harvard University?
	\begin{choices}
		\item To argue that happiness can be affected by various factors.
		\item To provide a historical context for the science of happiness.
		\item \textbf{\color{red}To support the view that happiness is significant to well-being.}
		\item To promote the idea that happiness is a pleasant emotional state.
	\end{choices}

	\item Why do some people fail to recognize happiness in everyday life?
	\begin{choices}
		\item \textbf{\color{red}They neglect to appreciate the small pleasures in life.}
		\item They lack meaningful connections with people around them.
		\item They focus too much on what they have experienced in the past.
		\item They don't know what responsibilities to take to achieve happiness.
	\end{choices}

	\item What does the speaker aim to convey about happiness?
	\begin{choices}
		\item Happiness is the key to leading a fulfilling and meaningful life.
		\item Happiness may not be a natural consequence of adult responsibilities.
		\item Happiness can be genuinely achieved after one realizes their goals in life.
		\item \textbf{\color{red}Happiness is a choice that can be made despite challenging circumstances.}
	\end{choices}
\end{enumerate}

\newpage
\subsection{Unit 6}

\subsubsection{Conversation}
\begin{enumerate}
	\item What is one of the reasons the man likes long novels?
	\begin{choices}
		\item The characters in long novels have complex problems to solve.
		\item The characters in long novels face unexpected challenges in their lives.
		\item The stories of long novels contain more innovative ideas than those of short novels.
		\item \textbf{\color{red}The stories of long novels are developed in greater depth than those of short novels.}
	\end{choices}

	\item What does the woman think of short novels?
	\begin{choices}
		\item They can hardly tell complex stories.
		\item They can be interpreted in a creative way.
		\item \textbf{\color{red}They can be as interesting as long novels.}
		\item They have room for plenty of fun twists.
	\end{choices}

	\item What did Hemingway do according to the conversation?
	\begin{choices}
		\item He condensed a long story into only several pages.
		\item \textbf{\color{red}He won a bet by writing a story in just several words.}
		\item He published a large number of mini stories in magazines.
		\item He wrote quite a few famous novels by recreating old legends.
	\end{choices}

	\item What do we learn about the American magazine's mini-story project mentioned in the conversation?
	\begin{choices}
		\item The project lasted for just several months.
		\item \textbf{\color{red}The project has attracted a lot of participants.}
		\item The project has received about 50,000 entries so far.
		\item The project inspires the participants to read more books.
	\end{choices}
\end{enumerate}

\subsubsection{Passage}
\begin{enumerate}
	\item What does the program mentioned in the passage do?
	\begin{choices}
		\item It sets the criteria for schoolchildren to pass reading exams.
		\item \textbf{\color{red}It helps schoolchildren overcome reading and writing problems.}
		\item It trains teachers to improve their current instructional methods.
		\item It studies the effects of literacy skills on students' success in school.
	\end{choices}

	\item What is a characteristic of Reading Recovery?
	\begin{choices}
		\item \textbf{\color{red}Lessons are taught in an interactive way.}
		\item Lessons last for an average period of 30 weeks.
		\item Students need to monitor each other's reading behavior.
		\item Students work in small groups to develop reading skills.
	\end{choices}

	\item When will students stop taking Reading Recovery lessons?
	\begin{choices}
		\item When they show good skills in assembling short stories.
		\item When they have developed a consistent reading behavior.
		\item \textbf{\color{red}When they have reached the average reading level for their grade.}
		\item When they demonstrate a high level of fluency in reading the alphabet.
	\end{choices}

	\item What is mentioned in the passage as a disadvantage of Reading Recovery?
	\begin{choices}
		\item The program is too short to have clear positive effects.
		\item The lesson materials can be too difficult for some students.
		\item The lessons demonstrate few long-term benefits for students.
		\item \textbf{\color{red}The costs of staff training are rather high for schools.}
	\end{choices}
\end{enumerate}

\subsubsection{Lecture 1}
\begin{enumerate}
	\item How can reading fiction affect our brains according to the speaker?
	\begin{choices}
		\item It can strengthen our memory.
		\item \textbf{\color{red}It can improve our neural networks.}
		\item It can reduce tension in some areas of the brain.
		\item It can increase our ability to express complex feelings.
	\end{choices}

	\item What is one of the activities we are engaged in while reading a novel according to the speaker?
	\begin{choices}
		\item \textbf{\color{red}Judging the meaning of literary texts.}
		\item Evaluating real-life situations creatively.
		\item Exploring links between complex stimuli.
		\item Analyzing the background of the narrator.
	\end{choices}

	\item What would the speaker most likely agree with regarding the effect of reading fiction?
	\begin{choices}
		\item Reading fiction may cause a decrease in our mental flexibility.
		\item Even reading just one fiction book brings clear cognitive changes.
		\item Reading fiction may result in unrealistic imagination of the world.
		\item \textbf{\color{red}Reading fiction for a long period can make lasting changes in the brain.}
	\end{choices}

	\item What does the passage mainly discuss?
	\begin{choices}
		\item Typical images and emotions created in fiction books.
		\item Mental skills commonly used to interpret literary texts.
		\item \textbf{\color{red}The cognitive and psychological benefits of reading novels.}
		\item Basic skills to develop empathy for others through reading.
	\end{choices}
\end{enumerate}

\subsubsection{Lecture 2}
\begin{enumerate}
	\item What is a characteristic of Victorian literature?
	\begin{choices}
		\item \textbf{\color{red}The novel was the dominant genre.}
		\item Styles created in previous eras were revived.
		\item Gold mining was a common topic in fictional works.
		\item Most writers were popular for just a short period of time.
	\end{choices}

	\item What do we know about novels in the Victorian period?
	\begin{choices}
		\item They were famous for advancing printing technology.
		\item They were crucial for improving the population's literacy.
		\item \textbf{\color{red}They were often printed in separate parts over a period of time.}
		\item They were often published in multiple formats to attract readers.
	\end{choices}

	\item What might a Victorian novelist be most interested in writing about?
	\begin{choices}
		\item Life changes in an abstract style.
		\item The bold new era of technology.
		\item \textbf{\color{red}Individual struggles in everyday life.}
		\item Romantic stories in a transforming society.
	\end{choices}

	\item What contributed to the prosperity of children's literature in the Victorian period?
	\begin{choices}
		\item The ever-increasing popularity of realistic literature.
		\item A commercial focus on literacy programs for children.
		\item \textbf{\color{red}Social programs focusing on protecting children's rights.}
		\item The rapid growth of a publishing market for young authors.
	\end{choices}
\end{enumerate}
